% ============================================================================
% Fichier  : partie0/P00_scene_crime_numerique.tex
% Titre    : Prologue -- Scène de Crime Numérique
% Auteur   : MINKA MI NGUIDJOÏ Thierry Emmanuel
% Version  : 1.0 -- Juin 2026
% Statut   : RÉDIGÉ
%
% Objectif pédagogique :
%   Créer l'urgence pédagogique dès la première page. L'étudiant comprend,
%   avant toute théorie, qu'une seule erreur procédurale annule des semaines
%   de travail technique. Le cas MBONGO sert de fil rouge immersif.
%   Les 4 principes fondamentaux sont posés ici dans leur forme la plus brute.
%   Le Trilemme CRO est introduit de façon opérationnelle.
% ============================================================================

\chapter{Prologue --- Scène de Crime Numérique}
\label{chap:prologue}

\epigraph{%
    « Un dossier judiciaire ne se perd jamais au tribunal.\\
    Il se perd sur la scène de crime, dans les dix premières minutes. »%
}{--- Adapté de la pratique forensique}

\niveauChapitre{Introductif}{%
    Aucun prérequis technique. Ce prologue se lit avant même d'ouvrir un
    terminal ou de brancher un disque dur.%
}

\begin{boxobjectifs}
\begin{itemize}
    \item Comprendre, à travers un cas concret, pourquoi la méthode prime
          sur la technique en investigation numérique.
    \item Identifier les quatre principes fondamentaux --- intégrité,
          traçabilité, reproductibilité, admissibilité --- et reconnaître
          une violation de chacun d'eux.
    \item Mémoriser les réflexes immédiats sur une scène de crime numérique
          (PC allumé, PC éteint, téléphone) avant toute manipulation.
    \item Saisir la nature du \termecle{Trilemme CRO} comme tension
          constitutive de toute investigation numérique.
\end{itemize}
\end{boxobjectifs}

\bigskip

% ============================================================================
\section*{Avant de commencer : une affaire qui n'aurait pas dû échouer}
% ============================================================================

Douala, mars~2026. Le tribunal de grande instance vient de rendre son verdict
dans l'affaire MBONGO\footnote{Les noms, dates et montants de ce cas sont
fictifs mais les mécanismes juridiques et techniques sont rigoureusement
exacts.}. Joseph MBONGO, comptable d'une PME de la Zone Industrielle de
Bassa, était accusé d'avoir détourné \textbf{12~millions de francs CFA} via
des virements frauduleux depuis le compte de son employeur vers un compte
tiers. Les preuves semblaient accablantes : un email de virement portant sa
signature, des captures d'écran du portail bancaire UBA sur son ordinateur,
un historique de navigation Chrome daté et horodaté, un fichier Word supprimé
nommé \texttt{faux\_bon\_de\_commande.docx}.

\medskip

L'avocat de la défense a demandé la parole. En dix minutes, il a obtenu
l'irrecevabilité de la totalité des preuves numériques. Le dossier s'est
effondré.

\begin{boxalert}
\textbf{Pourquoi ?} L'OPJ chargé de la saisie avait, en arrivant sur place,
\emph{allumé} l'ordinateur pour « constater visuellement les faits ». Windows
avait automatiquement monté toutes les partitions, mis à jour les horodatages
d'accès, créé des fichiers temporaires, modifié des clés de registre. L'avocat
a prouvé au juge que l'heure de « création » de la capture d'écran
\texttt{ecran\_virement.png}, telle que lue dans Autopsy, était
\textbf{postérieure} à l'heure de saisie déclarée dans le procès-verbal.
Contamination prouvée. Irrecevabilité prononcée.

L'OPJ avait pourtant fait \emph{tout le reste} correctement~: il avait utilisé
Autopsy, calculé les hash, rédigé un rapport de 40~pages. Tout cela est devenu
sans valeur à cause de \textbf{dix secondes d'inattention} au moment d'appuyer
sur le bouton Power.
\end{boxalert}

\medskip

Ce prologue n'est pas là pour vous enseigner Autopsy, ni le protocole TCP/IP,
ni la cryptographie post-quantique. Ces sujets occupent les 600~pages qui
suivent. Ce prologue est là pour vous enseigner \textbf{le réflexe fondateur
de toute la discipline}~: \emph{avant d'agir, comprendre ce que votre action
va détruire}.

\separateur

% ============================================================================
\section{La Scène de Crime Numérique : un Lieu à Sécuriser}
% ============================================================================

\subsection{Ce qui change par rapport à une scène physique}

Un investigateur criminaliste traditionnel sait qu'il ne doit pas marcher
sur les traces de pas, ni toucher l'arme sans gants. L'équivalent numérique
de cette règle est moins intuitif, parce que la contamination y est
\textbf{invisible}. Brancher un clavier, allumer un écran, ouvrir un
fichier~: chacune de ces actions laisse des traces qui modifient l'état
de la scène.

\begin{boxdef}
\textbf{Scène de crime numérique}\index{scène de crime numérique} --- Tout
environnement numérique susceptible de contenir des traces probantes d'un
acte délictueux. Elle inclut les supports physiques (disques durs, clés USB,
téléphones), la mémoire volatile (RAM), le trafic réseau, les journaux
de serveurs et les données stockées dans le nuage.
\end{boxdef}

La différence essentielle avec une scène physique est la
\termecle{volatilité}~: certaines données disparaissent en quelques
secondes (mémoire vive), d'autres en quelques minutes (connexions réseau),
d'autres subsistent mais peuvent être altérées sans laisser de trace
visible à l'~il nu (horodatages de fichiers).

\begin{tableaucours}{p{3.5cm}p{3.8cm}p{4.5cm}}{%
    \tableauHeader{Type de donnée} &
    \tableauHeader{Durée de vie} &
    \tableauHeader{Menace principale}
}
Registres CPU, cache & Nanosecondes & Non capturable \\
Mémoire RAM & Jusqu'à extinction & Extinction non planifiée \\
Connexions réseau actives & Secondes à minutes & Déconnexion réseau \\
Fichiers ouverts, sessions & Minutes & Fermeture de session \\
Disque dur / SSD & Mois à années & Chiffrement, TRIM \\
Journaux de serveur & Jours à semaines & Rotation automatique \\
Données cloud & Indéfini & Suppression distante \\
\end{tableaucours}
\vspace{-0.8em}
\begin{center}{\small\itshape Tableau~P.1 --- Ordre de volatilité des données numériques
(d'après RFC~3227 \cite{rfc3227})}\end{center}

\subsection{Le paradoxe de l'observation}

En physique quantique, le principe d'Heisenberg stipule que l'acte de mesurer
une particule modifie son état. L'investigation numérique connaît un paradoxe
analogue~: \textbf{tout acte d'observation d'un système numérique en cours
d'exécution le modifie}. Lancer un outil forensique sur un PC allumé écrit
dans la RAM, crée des fichiers temporaires, génère des entrées dans les
journaux système.

Ce paradoxe ne constitue pas une impasse --- il se gère par la
\textbf{documentation rigoureuse}~: l'investigateur note précisément chaque
action entreprise et chaque modification connue qu'elle a induite. Le tribunal
peut alors distinguer les traces \emph{antérieures} à l'investigation des
traces \emph{créées par} l'investigation.

\begin{boxinfo}
\textbf{La règle d'or~:} toute modification introduite pendant l'investigation
doit être \emph{documentée}, \emph{justifiée} et \emph{minimisée}. L'objectif
n'est pas la perfection impossible, mais la traçabilité complète.
\end{boxinfo}

% ============================================================================
\section{Les Quatre Principes Fondamentaux}
% ============================================================================

Ces quatre principes structurent l'intégralité de la discipline. Ils seront
formalisés mathématiquement dans la Partie~II, mais les voici dans leur
formulation la plus directe, illustrée par des violations réelles.

\subsection{Principe 1 --- Intégrité}
\label{subsec:principe-integrite}

\begin{boxdef}
\textbf{Intégrité}\index{intégrité!preuve numérique} (\textit{Integrity}) ---
Garantie que la preuve numérique n'a pas été modifiée depuis sa collecte.
Elle se démontre par la correspondance des \termecle{empreintes
cryptographiques} calculées avant et après acquisition.
\end{boxdef}

L'intégrité est prouvée \textbf{mathématiquement}, pas déclarativement.
Il ne suffit pas d'affirmer « je n'ai rien touché »~; il faut que le
\termecle{hash} de l'image forensique soit bit-à-bit identique au hash du
support original.

\begin{boxalert}
\textbf{Violation typique.} Un enquêteur branche un disque dur suspect
directement sur son ordinateur Windows sans \textit{write blocker}. Windows
monte la partition, met à jour les horodatages d'accès du système de fichiers
NTFS, crée un répertoire \texttt{System Volume Information}. Le hash du disque
après branchement diffère du hash avant branchement. La preuve est contaminée.

\smallskip
\textbf{Règle absolue~:} ne jamais monter un support suspect sans
\termecle{write blocker} matériel ou logiciel activé et vérifié.
\end{boxalert}

La vérification d'intégrité repose sur des fonctions de hachage
cryptographiques. Les commandes de référence~:

\begin{lstlisting}[language=bash_forensique,
    caption={Calcul et vérification d'empreintes cryptographiques}]
# Windows (CMD) -- calcul MD5 et SHA-256
certutil -hashfile C:\evidence\disque.E01 MD5
certutil -hashfile C:\evidence\disque.E01 SHA256

# Windows (PowerShell) -- alternative recommandée
Get-FileHash C:\evidence\disque.E01 -Algorithm SHA256

# Linux / Kali Linux -- syntaxe native
md5sum    /evidence/disque.E01
sha256sum /evidence/disque.E01
\end{lstlisting}

Les deux hash --- source et image --- doivent être \textbf{identiques} et
\textbf{documentés} dans le formulaire de chaîne de custody \emph{avant}
toute analyse.

\begin{boiteRemarque}{MD5 ou SHA-256 ?}
MD5 (128~bits) reste accepté en forensique pour l'intégrité d'une image
car le risque n'est pas la falsification intentionnelle mais l'altération
accidentelle. SHA-256 (256~bits) est recommandé pour toutes les nouvelles
investigations et exigé en contexte à haute valeur probatoire. Calculez
\emph{toujours les deux simultanément} et documentez les deux dans le
formulaire de custody.
\end{boiteRemarque}

\subsection{Principe 2 --- Traçabilité (chaîne de custody)}
\label{subsec:principe-tracabilite}

\begin{boxdef}
\termecle{Chaîne de custody}\index{chaîne de custody} (\textit{Chain of
custody}) --- Registre chronologique documentant chaque personne ayant eu
accès à une preuve et toutes les actions effectuées sur celle-ci, depuis
sa collecte jusqu'à sa présentation en audience.
\end{boxdef}

La chaîne de custody est le \textbf{document le plus important d'une
investigation numérique} --- plus important que le rapport technique
lui-même. Un rapport d'analyse de 100~pages perd toute valeur si on ne peut
prouver que l'evidence analysée est bien celle saisie sur la scène de crime.

\begin{boxcustody}
\textbf{Le formulaire commence au premier contact.}
Les informations minimales à consigner immédiatement~:
\begin{enumerate}
    \item Numéro de dossier et numéro d'evidence (\textit{ex.}~: DOS-2026-004 / EV-001)
    \item Description physique précise~: marque, modèle, numéro de série,
          couleur, état apparent (rayures, caches manquants)
    \item Lieu, date et heure de collecte (\texttt{JJ/MM/AAAA HH:MM})
    \item Nom, grade et matricule du collecteur
    \item \hash{MD5}{} et \hash{SHA-256}{} du support original (avant acquisition)
    \item \hash{MD5}{} et \hash{SHA-256}{} de l'image forensique (après acquisition)
    \item Confirmation explicite que les deux hash correspondent
\end{enumerate}
Le formulaire complet est disponible en Annexe~C.
\end{boxcustody}

\begin{boxalert}
\textbf{Violation typique.} Un OPJ saisit le téléphone d'un suspect et le
confie à un collègue pendant une heure pour « garder un \oe il dessus ».
Ce transfert non documenté constitue un vide dans la chaîne de custody.
À l'audience, l'avocat peut semer le doute~: le téléphone a-t-il été
manipulé pendant cette heure ?

\textbf{Règle~:} tout transfert d'une evidence entre deux personnes doit
être documenté immédiatement --- date, heure, raison, signatures des deux
parties. Si un doute existe, l'intégrité doit être re-vérifiée par re-hashage.
\end{boxalert}

\begin{boxjuris}
\textbf{Droit camerounais --- \loicam{2010/012}, article~53, alinéa~1~:}

\medskip
\textit{« Les perquisitions \textelp{} peuvent porter sur des copies réalisées
en présence des personnes qui assistent à la perquisition. »}

\medskip
Cet article établit l'exigence légale d'une copie forensique réalisée
\emph{en présence} des parties. Le formulaire de chaîne de custody est la
traduction procédurale de cette exigence~: il prouve que la copie a été
réalisée de façon transparente et documentée. Son absence peut suffire à
faire rejeter l'ensemble du dossier numérique.
\end{boxjuris}

\subsection{Principe 3 --- Reproductibilité}
\label{subsec:principe-reproductibilite}

\begin{boxdef}
\textbf{Reproductibilité}\index{reproductibilité} (\textit{Reproducibility}) ---
Propriété d'une méthode d'investigation telle que, si elle est appliquée
par un autre investigateur qualifié sur la même evidence avec les mêmes
outils, elle produit les mêmes résultats.
\end{boxdef}

La reproductibilité est la \textbf{condition de scientificité} de
l'investigation numérique. Un résultat que seul son auteur peut obtenir,
avec un outil non documenté, selon une méthode non explicitée, n'est pas
une preuve scientifique~: c'est une opinion.

En pratique, trois exigences concrètes s'imposent~:

\begin{enumerate}
    \item \textbf{Nommer les outils avec précision~:} nom, version exacte,
          système d'exploitation hôte. \textit{« Analysé avec Autopsy »}
          est insuffisant. La formulation correcte est~: \textit{« Autopsy
          4.21.0 (The Sleuth Kit 4.12.0) sur Windows~11 Pro~22H2. »}

    \item \textbf{Décrire la méthode~:} les étapes dans leur ordre exact,
          suffisamment détaillées pour qu'un pair puisse les reproduire
          indépendamment.

    \item \textbf{Documenter les limites~:} ce qui n'a \emph{pas} pu être
          analysé (partition chiffrée, fichier corrompu) doit être
          explicitement mentionné. Un rapport silencieux sur ses propres
          lacunes est un rapport trompeur.
\end{enumerate}

\begin{boxalert}
\textbf{Violation typique.} Un analyste utilise un script Python maison
pour extraire des artefacts de navigation. L'expert adverse tente de
reproduire l'analyse sur la même image --- et obtient des résultats
différents, car le script contient une erreur sur certains formats SQLite.
La méthode n'est pas reproductible. Les conclusions sont invalidées.

\textbf{Règle~:} utiliser des outils validés et documentés (Autopsy,
Volatility, FTK Imager). Documenter la version exacte. Tester tout outil
maison sur des données de référence connues avant utilisation en investigation
réelle.
\end{boxalert}

\subsection{Principe 4 --- Admissibilité}
\label{subsec:principe-admissibilite}

\begin{boxdef}
\textbf{Admissibilité}\index{admissibilité} (\textit{Admissibility}) ---
Qualité d'une preuve numérique lui permettant d'être recevable dans une
procédure judiciaire. Elle suppose que la preuve a été obtenue légalement,
qu'elle est pertinente aux faits allégués, et qu'elle est intègre.
\end{boxdef}

Une preuve techniquement parfaite est sans valeur si elle est juridiquement
irrecevable. L'admissibilité est la \emph{finalité} de tous les autres
principes~: l'intégrité, la traçabilité et la reproductibilité existent
\emph{pour} garantir l'admissibilité.

\begin{boxalert}
\textbf{Violation typique.} Un supérieur hiérarchique accède au poste de
travail d'un employé suspecté avec son propre compte administrateur, copie
les fichiers sur une clé USB personnelle et les remet à la police. Ces
fichiers sont irrecevables~: collecte sans autorisation judiciaire,
procédure de saisie non conforme aux articles~52--59 de la
\loicam{2010/012}.

\textbf{Règle~:} en contexte judiciaire, l'investigation numérique exige
l'autorisation appropriée (mandat de perquisition, réquisition du procureur).
En contexte interne (entreprise), elle doit être encadrée par le règlement
intérieur et respecter le droit du travail camerounais.
\end{boxalert}

\begin{boxjuris}
\textbf{Quatre critères d'admissibilité en droit camerounais}\\
\textit{(d'après la jurisprudence et les articles~52--59 de la \loicam{2010/012})}

\begin{enumerate}
    \item \textbf{Légalité~:} obtenue avec l'autorisation requise par la loi.
    \item \textbf{Pertinence~:} en lien direct avec les faits poursuivis.
    \item \textbf{Intégrité~:} non altérée depuis sa collecte (hachage vérifiable).
    \item \textbf{Authenticité~:} origine prouvable par la chaîne de custody.
\end{enumerate}
\end{boxjuris}

% ============================================================================
\section{Le Trilemme CRO : la Tension Constitutive}
\label{sec:trilemme-intro}
% ============================================================================

Avant d'entrer dans le détail des techniques, il est essentiel de comprendre
la tension profonde qui traverse toute l'investigation numérique. Cette tension
a été formalisée par les travaux du \sigle{LIMSI} sous le nom de
\termecle{Trilemme CRO}\index{Trilemme CRO} \cite{minka2025cro}.

\subsection{Les trois dimensions en opposition}

Toute investigation numérique doit simultanément satisfaire trois exigences
qui se contrarient mutuellement~:

\begin{itemize}
    \item \textbf{$\mathcal{C}$ --- Confidentialité~:} protéger la vie privée
          des personnes impliquées (suspect, victimes, témoins) et les données
          sensibles rencontrées.
    \item \textbf{$\mathcal{R}$ --- Fiabilité~(\textit{Reliability})~:} garantir
          que les preuves sont intègres, authentiques, reproductibles.
    \item \textbf{$\mathcal{O}$ --- Opposabilité juridique~:} assurer que les
          preuves peuvent être présentées et défendues devant un tribunal.
\end{itemize}

\begin{boxCRO}
\textbf{Le Trilemme formulé~:} il est impossible de maximiser simultanément
les trois dimensions.

\medskip
Pour qu'une preuve soit \textbf{opposable} ($\mathcal{O}$ haute), elle doit
être vérifiable par toutes les parties --- ce qui tend à \emph{réduire} la
confidentialité ($\mathcal{C}$ contrainte).

Pour maximiser la \textbf{confidentialité} ($\mathcal{C}$ haute) --- en
chiffrant les données collectées --- on risque de compromettre la
reproductibilité et donc la fiabilité ($\mathcal{R}$ réduite).

Pour garantir la \textbf{fiabilité} ($\mathcal{R}$ haute) --- en collectant
le maximum de données --- on élargit le périmètre et empiète sur la vie
privée de personnes non impliquées ($\mathcal{C}$ réduite).

\[
\Gamma_{CRO}(\Pi) \;=\; \max\bigl\{\mathcal{C}(\Pi),\;
\mathcal{R}(\Pi),\; \mathcal{O}(\Pi)\bigr\}
\;\geq\; 0.4 + \mathrm{negl}(\lambda)
\]

\textit{Toute investigation est un compromis explicite entre ces trois
dimensions. La maîtrise de ce compromis distingue l'investigateur expert
du technicien.}
\end{boxCRO}

\subsection{Le Trilemme sur la scène de crime : exemples immédiats}

Ce n'est pas une abstraction théorique. Dès les premières secondes,
l'investigateur fait face à des choix CRO concrets~:

\begin{itemize}
    \item \textit{Doit-il capturer la RAM ?} Cette capture va légèrement
          modifier la RAM (atteinte à $\mathcal{R}$, documentable et
          acceptable) mais c'est la seule façon d'obtenir les clés de
          chiffrement actives (gain majeur pour $\mathcal{O}$).
          \hfill\cro{$\sim$}{$\downarrow$faible}{$\uparrow$}

    \item \textit{Peut-il lire tous les messages WhatsApp ?} Certains n'ont
          aucun lien avec l'affaire et révèlent la vie privée de tiers
          innocents. Lire plus que nécessaire améliore la fiabilité brute
          mais viole la proportionnalité légale.
          \hfill\cro{$\downarrow$}{$\uparrow$}{$\sim$}

    \item \textit{Doit-il tout documenter publiquement ?} Certains artefacts
          sensibles (données médicales, correspondances protégées) peuvent
          être transmis sous pli scellé au juge. Gain pour $\mathcal{C}$,
          légère contrainte pour $\mathcal{O}$.
          \hfill\cro{$\uparrow$}{$\sim$}{$\downarrow$faible}
\end{itemize}

\medskip

Tout au long de ce cours, chaque technique sera analysée à travers ce prisme.
L'encadré \textbf{Analyse~CRO} apparaîtra systématiquement pour expliciter
quel compromis est réalisé et comment le justifier devant un tribunal.

% ============================================================================
\section{Protocole d'Urgence : les 10 Premières Minutes}
\label{sec:protocole-urgence}
% ============================================================================

Ce protocole synthétise les décisions critiques des dix premières minutes.
Il sera détaillé et justifié dans la Partie~II et la Partie~IV. Il est
présenté ici dans sa forme la plus opérationnelle, pour être applicable
immédiatement.

\subsection{Arrivée sur la scène : les trois questions immédiates}

\begin{enumerate}
    \item \textbf{Le matériel est-il sous tension ?} → La réponse détermine
          tout ce qui suit.
    \item \textbf{Y a-t-il du chiffrement actif visible ?}
          (BitLocker déverrouillé, VeraCrypt monté, Android déverrouillé) →
          Éteindre un support chiffré sans le capturer d'abord peut rendre
          les données définitivement inaccessibles.
    \item \textbf{Y a-t-il une connexion réseau active ?} → Un suspect peut
          encore déclencher une destruction de données à distance tant que
          la connexion est maintenue.
\end{enumerate}

\subsection{PC allumé : procédure de triage}

\begin{boxPNC}
\textbf{PC allumé --- ordre des actions}

\begin{enumerate}
    \item \volatilite{1}{RAM} --- Capturer la mémoire vive avec FTK~Imager
          (\texttt{File → Capture Memory}). Le PC reste allumé pendant
          cette étape.
    \item \volatilite{2}{Écran} --- Photographier l'écran dans son état
          actuel. Ne fermer aucune fenêtre.
    \item \volatilite{3}{Réseau} --- Documenter les connexions actives~:
          \texttt{netstat -ano} (CMD) ou \texttt{Get-NetTCPConnection}
          (PowerShell). Photographier ou noter.
    \item \volatilite{4}{Arrêt} --- Éteindre proprement via le menu Windows.
          Sauf si une destruction de données est imminente, auquel cas
          l'arrachage du câble est justifiable et \emph{doit être documenté}.
    \item \volatilite{5}{Disque} --- Brancher avec write blocker. Acquérir
          l'image E01 avec FTK~Imager. Calculer et noter hash MD5+SHA-256.
\end{enumerate}

\medskip
\textbf{Ne jamais~:} ouvrir un fichier, lancer une application, connecter
un périphérique USB sans write blocker, laisser Windows exécuter une tâche
automatique (mise à jour, indexation).
\end{boxPNC}

\subsection{PC éteint : procédure de saisie}

\begin{boxPNC}
\textbf{PC éteint --- ordre des actions}

\begin{enumerate}
    \item \textbf{Ne pas allumer.} Appuyer sur le bouton Power contamine
          la scène avant même que l'OS ne soit chargé.
    \item Photographier le matériel sous tous les angles, câbles en place,
          avant toute manipulation physique.
    \item Activer le write blocker logiciel ou utiliser un write blocker
          matériel, puis retirer le disque dur ou SSD.
    \item Acquérir l'image forensique en format E01 avec FTK~Imager.
          Documenter les métadonnées du cas (numéro de dossier, examinateur).
    \item Calculer et documenter les hash MD5 et SHA-256 source et image.
    \item Sceller l'evidence originale dans un sachet inviolable étiqueté
          et consigner dans le formulaire de chaîne de custody.
\end{enumerate}
\end{boxPNC}

\subsection{Téléphone mobile : les règles absolues}

\begin{boxPNC}
\textbf{Téléphone mobile --- La règle d'or~: mode avion en premier.}

\medskip
Un téléphone connecté peut recevoir une commande d'effacement à distance
(\textit{remote wipe}) en quelques secondes. Le mode avion coupe la
connexion sans effacer les données.

\medskip
\begin{enumerate}
    \item \textbf{Téléphone allumé~:} mode avion immédiat. Si l'écran est
          inaccessible, placer dans un sac de Faraday\footnote{À défaut de
          sac de Faraday~: plusieurs épaisseurs d'aluminium ménager
          hermétiquement scellées. Vérifier que le signal opérateur disparaît.}.
          Ne jamais brancher sur un chargeur avant isolation électromagnétique.

    \item \textbf{Téléphone éteint, Android, code connu~:} allumer en mode
          avion (maintenir le bouton Power et activer le mode avion dès
          l'affichage du démarrage). Extraire via ADB.

    \item \textbf{Téléphone éteint, chiffré, code inconnu~:} conserver
          éteint. Alimenter la batterie si nécessaire via câble contrôlé.
          Transmettre à un laboratoire spécialisé. N'allumer en aucun cas
          sans stratégie d'extraction préparée au préalable.

    \item \textbf{Documenter~:} numéro IMEI (\texttt{*\#06\#} ou sous la
          batterie), opérateur, état de l'écran, niveau de batterie estimé.
\end{enumerate}
\end{boxPNC}

% ============================================================================
\section{Retour sur l'Affaire MBONGO : ce qui Aurait Dû se Passer}
\label{sec:retour-mbongo}
% ============================================================================

Voici la reconstitution de l'affaire d'ouverture, conduite cette fois
conformément aux quatre principes.

\begin{boxscenario}{Affaire MBONGO --- Version correcte}
\textbf{Situation initiale~:} PC Windows allumé et déverrouillé.
Téléphone Android sur le bureau, écran actif. L'OPJ arrive sur place avec
une clé USB write-protégée contenant FTK~Imager Portable.

\medskip
\begin{enumerate}
    \item \textbf{Photographies d'état} --- L'OPJ photographie l'écran PC
          (fenêtre navigateur ouverte sur le portail bancaire UBA,
          horodatage visible dans la barre des tâches) et l'écran du téléphone.
          Il ne touche ni l'un ni l'autre.

    \item \textbf{Mode avion} sur le téléphone. Consigné dans le formulaire
          (\textit{« 14h37~: téléphone passé en mode avion par OPJ NKENG »}).

    \item \textbf{Capture RAM} --- FTK~Imager Portable lancé depuis la clé
          USB. \texttt{File → Capture Memory}. Fichier \texttt{RAM\_MBONGO.mem}
          sauvegardé sur une clé de destination séparée. Taille~: 16~Go.
          Hash MD5 calculé et noté : \hash{MD5}{a3f7...}.

    \item \textbf{Connexions réseau} --- \texttt{netstat -ano} lancé dans
          CMD. Résultat photographié~: deux connexions ESTABLISHED vers
          l'IP \texttt{41.202.219.X} (UBA~Cameroun) et une connexion vers
          \texttt{185.62.X.X} (inconnue, notée pour investigation ultérieure).

    \item \textbf{Arrêt propre} du PC via le menu Démarrer. Retrait du
          disque dur après déconnexion. Write blocker logiciel activé.
          Acquisition E01 en 47~minutes. Hash source~: \hash{SHA-256}{3d9c...}.
          Hash image~: \hash{SHA-256}{3d9c...}. \textbf{Intégrité vérifiée.}

    \item \textbf{Formulaire de chaîne de custody} rempli pour le PC et le
          téléphone. Sachets scellés. Deux signatures. Transport au dépôt
          de preuves de la Brigade de Cybercriminalité.

    \item \textbf{Analyse sur copie} --- Autopsy~4.21.0 ouvre l'image E01.
          Jamais le disque original. Les 8~artefacts sont identifiés~;
          leurs horodatages sont tous \emph{antérieurs} à la saisie.
\end{enumerate}

\textbf{Résultat~:} à l'audience, l'expert présente les hash correspondants,
le formulaire de custody ininterrompu, la version exacte des outils. L'avocat
de la défense ne peut pas arguer d'une contamination postérieure à la saisie.
Le dossier tient.
\end{boxscenario}

\begin{boxCRO}
\textbf{Analyse CRO --- Affaire MBONGO}

\begin{itemize}
    \item $\mathcal{C}$ --- La capture RAM révèle peut-être des données sans
          rapport avec l'affaire. Ce risque est mentionné dans le rapport~;
          les données non pertinentes ne sont pas incluses dans le dossier.
          \hfill\textit{Compromise acceptable et documentée.}

    \item $\mathcal{R}$ --- La capture RAM modifie légèrement la RAM
          elle-même (le processus FTK y réside pendant l'acquisition).
          Modification connue, quantifiable, minimale. Hash pris
          \emph{après} capture pour refléter cet état.
          \hfill\textit{Maîtrisée.}

    \item $\mathcal{O}$ --- Le hachage, le formulaire de custody et les
          versions d'outils documentés permettent à l'expert de défendre
          chaque artefact devant le juge.
          \hfill\textit{Maximisée.}
\end{itemize}

\cro{$\downarrow$\,acceptable}{$\sim$\,maîtrisé}{$\uparrow$\,maximisé} ---
Configuration standard d'une investigation judiciaire correctement conduite.
\end{boxCRO}

% ============================================================================
\separateur
% ============================================================================

\begin{boxresume}
\textbf{Ce qu'il faut retenir de ce prologue~:}

\begin{itemize}
    \item Une investigation numérique peut être techniquement irréprochable
          et juridiquement anéantie par une erreur procédurale commise dans
          les premières minutes.

    \item Les \textbf{quatre principes fondamentaux} structurent toute la
          discipline~: \textbf{intégrité} (hachage MD5+SHA-256),
          \textbf{traçabilité} (chaîne de custody continue),
          \textbf{reproductibilité} (méthodes et outils documentés),
          \textbf{admissibilité} (légalité + pertinence + intégrité).

    \item L'intégrité se prouve \textbf{mathématiquement}, jamais
          déclarativement. Deux hash identiques (source et image) constituent
          la preuve.

    \item La \textbf{chaîne de custody} commence à la seconde où
          l'investigateur touche l'evidence --- pas quand il rédige son rapport.

    \item Le \termecle{Trilemme CRO} formalise la tension irréductible entre
          confidentialité~($\mathcal{C}$), fiabilité~($\mathcal{R}$)
          et opposabilité~($\mathcal{O}$). Toute investigation est un
          compromis explicite et documenté entre ces trois dimensions.

    \item Sur une scène numérique~: \textbf{observer} avant d'agir,
          \textbf{agir dans l'ordre de volatilité}, \textbf{documenter
          chaque geste}.
\end{itemize}
\end{boxresume}

\bigskip

\begin{boxexo}[title={\ding{45}\quad Exercice P.1 --- Diagnostic d'une violation \niveaubase}]
Lisez la description ci-dessous. Identifiez~: (a)~quel(s) principe(s)
fondamentaux ont été violés~; (b)~quelle dimension du Trilemme CRO a été
sacrifiée~; (c)~ce qu'il aurait fallu faire.

\medskip
\textit{Un chef d'entreprise suspecte un employé de fuite de données
vers un concurrent. Sans prévenir les autorités, il demande à son
informaticien de copier le disque dur de l'employé sur une clé USB
personnelle et de « fouiller les fichiers ». L'informaticien ouvre
les fichiers directement, prend des captures d'écran, et remet la
clé au chef d'entreprise. Trois semaines plus tard, l'entreprise
dépose plainte et présente ces éléments au procureur.}
\end{boxexo}

\begin{boxexo}[title={\ding{45}\quad Exercice P.2 --- Scénario de triage \niveaumoyen}]
Vous arrivez sur une scène avec~:
\begin{itemize}
    \item Un ordinateur portable Windows~11, allumé, écran déverrouillé,
          fenêtre de navigateur ouverte sur un portail bancaire.
    \item Un téléphone Android, allumé, sans code PIN visible, Wi-Fi actif.
    \item Une clé USB branchée sur le laptop, voyant clignotant.
\end{itemize}
Décrivez, dans l'ordre exact des priorités, les six premières actions que
vous effectuez. Pour chaque action, précisez~: le principe fondamental
qu'elle préserve, l'ordre de volatilité auquel elle répond, et l'impact
CRO (\cro{?}{?}{?}).
\end{boxexo}

\begin{boxexo}[title={\ding{45}\quad Exercice P.3 --- Réflexion CRO \niveauexpert}]
Dans l'affaire MBONGO, l'OPJ a capturé les 16~Go de RAM du PC. Cette RAM
contenait~: les preuves des virements, les identifiants de session bancaire,
\emph{et} 3~Go de correspondances personnelles de MBONGO sans rapport avec
l'affaire (échanges WhatsApp~Web, photos personnelles).

L'avocat de la défense plaide que ces données personnelles n'auraient pas
dû être collectées (violation de la vie privée, article~17 du Pacte
international relatif aux droits civils et politiques, applicable au
Cameroun depuis 1984).

\smallskip
(a)~Comment l'investigateur aurait-il dû anticiper ce problème avant
la capture~?\\
(b)~Comment le gérer dans son rapport pour préserver à la fois
$\mathcal{O}$ et $\mathcal{C}$~?\\
(c)~Existe-t-il une solution technique qui aurait permis de réduire
la surface de collecte~?
\end{boxexo}

\bigskip

\begin{boxinfo}
\textbf{La suite du cours.} Ce prologue vous a donné les réflexes de base
et la tension fondamentale (Trilemme CRO) qui donne sens à toutes les
techniques. La \textbf{Partie~I} (Chapitres~1--5) apporte la philosophie
et l'histoire qui ancrent ces réflexes dans leur contexte intellectuel.
La \textbf{Partie~II} (Chapitres~6--9) formalise le Processus Investigatif
Universel, la chaîne de custody et les outils d'acquisition. La
\textbf{Partie~IV} (Chapitres~14--20) vous entraîne sur des systèmes réels
--- Windows, réseau, mobile, cloud --- avec les cas qui traversent ce cours.
\end{boxinfo}
